% !TeX program = xelatex
\documentclass[10.5pt,compsoc,UTF8]{CjC}
\usepackage{CTEX}
\usepackage{graphicx}
\usepackage{footmisc}
\usepackage{multirow}
\usepackage{multicol}
\usepackage{amsmath,amsthm}
\usepackage{amssymb,amsfonts}
\usepackage{booktabs}
\usepackage{fancyhdr}
\usepackage{caption}
\usepackage{xcolor}
\usepackage{captionhack}
\usepackage{gbt7714}
\usepackage{array}
\usepackage{float}
\usepackage{needspace}
\usepackage{etoolbox}
\setcounter{page}{1}
\graphicspath{{figures/}}

%=======页眉设置=======%
\headevenname{\mbox{\quad} \hfill \mbox{\zihao{-5}{计\quad 算\quad 机\quad 学\quad 报} \hspace{50mm} \mbox{2026 年}}}%
\headoddname{第 1 卷 第 1 期 \hfill 大语言模型检索增强生成技术综述}
%=======页眉设置=======%

\renewcommand{\thefootnote}{\fnsymbol{footnote}}
\setcounter{footnote}{0}
\renewcommand\footnotelayout{\zihao{5-}}
\newtheoremstyle{mystyle}{0pt}{0pt}{\normalfont}{1em}{\bf}{}{1em}{}
\theoremstyle{mystyle}
\renewcommand\figurename{figure~}
\newcommand{\upcite}[1]{\textsuperscript{\cite{#1}}}
\renewcommand{\labelenumi}{(\arabic{enumi})}
\setlength\parindent{2em}
% 分栏断点：标题不落在栏底
% \pretocmd{\section}{\needspace{5\baselineskip}}{}{}
% \pretocmd{\subsection}{\needspace{3\baselineskip}}{}{}
\setlength\emergencystretch{1.5em}

\pagestyle{CjCheadings}
\IfFontExistsTF{Times New Roman}{\setmainfont{Times New Roman}}{\setmainfont{Liberation Serif}}

%===========================================================%
\begin{document}

\hyphenpenalty=50000
\makeatletter
\newcommand\mysmall{\@setfontsize\mysmall{7}{9.5}}
\makeatother
\let\temp\footnote
\renewcommand \footnote[1]{\temp{\zihao{-5}#1}}

%=======首页页眉=======%
\thispagestyle{empty}%
% \twocolumn[{%
\vspace{-13mm}
\noindent\begin{tabular}{p{168mm}}
\zihao{5-}
第 1 卷\quad 第 1 期 \hfill 计\quad 算\quad 机\quad 学\quad 报\hfill Vol. 1\quad No. 1\\
\zihao{5-}
2026年11月 \hfill CHINESE JOURNAL OF COMPUTERS \hfill Nov. 2026\\
\hline\\[-5.5mm]
\hline\end{tabular}
%=======首页页眉=======%
\par
%中文标题、作者与注脚
{
\centering
\vspace {11mm}
{\zihao{2} \heiti 大语言模型检索增强生成技术综述：\\[2mm]从朴素检索到智能体化与图结构增强 }
\par
\vskip 5mm
\par
{\zihao{3}\fangsong 陈庠宇$^{1)}$}%
\par\vspace {2mm}%
\zihao{6}{$^{1)}$(华侨大学)}\par
\vspace {2mm}%
{\zihao{6} 本报告为《专业性应用研究方法与运用》课程作业；作者为硕士研究生，主要研究领域为大语言模型、信息检索与检索增强生成。E-mail: 作者邮箱@example.com}\par
}%中文标题、作者
\par
\vskip 5mm
\par
\zihao{5-}{
\setlength{\baselineskip}{16pt}\selectfont{
\noindent {\heiti 摘\quad 要\quad }
大语言模型在开放域问答等知识密集型任务上表现突出，但其参数化知识存在时效滞后、长尾覆盖不足以及生成结论难以溯源的问题。检索增强生成（retrieval-augmented generation，RAG）利用外部知识源，如文档、数据库或结构化数据，提升模型性能，生成更准确、更符合上下文的输出。本调查以技术演进为主线，将检索增强生成的研究划分为朴素检索增强、进阶检索增强、模块化检索增强以及智能体化与图结构增强四个阶段，逐一梳理各阶段的核心机制、代表性方法与适用条件；在此基础上从基准数据集、评价指标与实证结论三个层面归纳现有评测体系，并讨论"检索是否总是有益"这一边界问题；最后总结领域共识、证据分歧与尚存不足，提出面向自适应检索、端到端联合优化、结构化与多模态知识接入以及可溯源性保障的研究方向。分析表明，检索增强生成效果并不取决于检索模块的复杂度，而取决于检索时机、证据粒度与生成器利用证据的方式三者之间的匹配程度。\par}}
\par
\vspace {5mm}
\par
\zihao{5-}{\noindent
{\heiti 关键词 \quad }{大语言模型；检索增强生成；稠密检索；知识密集型任务}
}\par\noindent
\par
\vskip 5mm
% }]%

\begin{multicols}{2}
\section{\heiti 引言}

\subsection{\heiti 研究背景与动机}

以 Transformer 为骨架的大语言模型通过大规模预训练将海量世界知识压缩进参数，从而在开放域问答、摘要与推理等任务上取得了显著进展。然而，这种"知识全部写在参数里"的范式存在三个结构性缺陷：其一，参数化知识的更新依赖于重新训练或继续预训练，成本高昂且难以跟上知识的时效变化；其二，长尾、专业与私有领域的知识在预训练语料中占比有限，模型容易给出貌似合理却缺乏依据的答案，即所谓幻觉；其三，模型输出缺少可核验的证据链，使用者难以判断某一句结论来自何处。检索增强生成（RAG）正是针对这三点提出的：它在推理阶段从外部语料中检索相关片段，并把片段作为条件输入生成器，使模型"带着从外部更新的知识作答"，同时把知识更新转化为索引更新\upcite{lewis2020retrieval,gao2023retrieval}。

近五年的研究使 RAG 从单一的"检索—拼接—生成"结构演化为包含索引构建、查询改写、多路召回、重排、上下文压缩、迭代检索与工具调用的复杂系统。工程实践的推进速度超过了概念梳理的速度：一方面，不同工作对"检索增强"的定义宽严不一，有的指在预训练阶段引入检索器\upcite{guu2020realm,borgeaud2021improving}，有的指在推理阶段拼接上下文\upcite{ram2023context}；另一方面，评测口径分散，同一方法在不同基准上的结论常常相反，导致"检索是否有效"这一问题至今缺少共识性回答\upcite{li2024retrieval,liu2024lost}。因此，以技术演进为线索对已有工作进行结构化梳理，既有助于厘清方法之间的继承关系，也有助于识别被实证结果支持与被证据反驳的论断。

\subsection{\heiti 研究范围与核心问题}

本文将讨论范围限定为面向自然语言知识密集型任务的检索增强生成，即生成器的输入除用户查询外还包含从外部语料检索到的证据，不包含仅依赖工具调用而无关知识检索的智能体工作，也不展开多模态检索的模型细节。围绕这一范围，本文试图回答四个问题：

\begin{enumerate}
\item 检索增强生成经历了怎样的技术演进，各阶段的划分依据是什么；
\item 各阶段的核心机制如何缓解检索与生成之间的失配，其代价与适用条件分别是什么；
\item 在什么条件下检索真正带来效果，现有实证证据在哪些问题上一致、在哪些问题上冲突；
\item 尚未解决的关键问题是什么，哪些方向具备可验证的研究价值。
\end{enumerate}

就研究意义而言，RAG 是当前把大模型接入动态、私有与专业知识的成本最低的路径之一。厘清其能力边界，既能避免在检索明显无益的场景中投入工程成本，也能为面向可信生成的系统设计提供依据。

\subsection{\heiti 文献检索与筛选}

本文以 DBLP、arXiv、ACL Anthology、Google Scholar 与 Web of Science 为检索源，检索式由主题词与同义词组合而成，英文包括 retrieval-augmented generation、retrieval augmentation、dense retrieval、knowledge-intensive NLP，中文包括检索增强生成、稠密检索、知识密集型任务。检索时间跨度设为 2020 年 1 月至 2025 年 3 月，并对 2020 年以前关于开放域问答与最近邻语言模型的奠基性工作做追溯检索。

本文的纳入标准为：(1) 提出或系统评述检索增强流程中至少一个环节的方法；(2) 给出可复现的实验或公开的基准数据；(3) 提供可检索的完整出处。排除标准为：(1) 仅讨论工程实现而无方法贡献的技术博客与产品文档；(2) 无实验支撑的立场性短文；(3) 与知识检索无关的纯提示工程工作。经去重与相关性筛选，最终纳入 60 篇文献，其中 2023 年之后发表的占多数，兼顾了奠基性工作与近期进展。

\subsection{\heiti 本文组织结构}

全文共分六节。第 1 节交代研究背景、研究范围与核心问题，并说明文献检索与筛选过程；第 2 节给出检索增强生成的问题形式化与标准流水线，为后续讨论提供统一术语；第 3 节按朴素检索增强、进阶检索增强、模块化检索增强以及智能体化与图结构增强四个阶段梳理技术演进，并对代表性方法做横向比较；第 4 节从基准数据集、评价指标与实证结论三个层面讨论评测体系，归纳领域共识与现有分歧；第 5 节总结尚存问题并给出未来研究方向；第 6 节概括全文结论。

\section{\heiti 检索增强生成的基础框架}

\subsection{\heiti 问题形式化}

设查询为 $x$、答案为 $y$、检索到的证据片段为 $z$。检索增强生成把答案的分布写成对证据的边缘化形式\upcite{lewis2020retrieval}：

\begin{equation}
\begin{aligned}
p(y\mid x)&=\sum_{z\in\mathcal{Z}} p_{\eta}(z\mid x)\,p_{\theta}(y\mid x,z)\\
&\approx\sum_{z\in\mathrm{TopK}(p_{\eta}(\cdot\mid x))} p_{\eta}(z\mid x)\,p_{\theta}(y\mid x,z),
\end{aligned}
\label{eq:rag}
\end{equation}

其中 $p_{\eta}(z\mid x)$ 为检索器（参数 $\eta$），$p_{\theta}(y\mid x,z)$ 为生成器（参数 $\theta$），$\mathcal{Z}$ 为整个语料上的片段集合。由于对全语料求和不可行，实践中用检索器返回的前 $k$ 个片段近似，这一近似既是效率的来源，也是误差的来源：当正确证据未被召回时，生成器只能在错误证据上作答，误差被后续环节放大\upcite{chen2024benchmarking}。

按证据参与方式的不同，检索增强可分为两类极端：一类把检索结果拼接进输入上下文，生成器结构不变，因而适用于黑盒模型\upcite{ram2023context,shi2023replug}；另一类在生成器内部以交叉注意力或隐变量方式融合证据，需要联合训练\upcite{lewis2020retrieval,izacard2020leveraging}。前者部署成本低但受上下文窗口限制，后者效果上限更高但训练开销大，二者的取舍贯穿后续各阶段的演化。

\subsection{\heiti 检索增强的标准工作流程}

一个完整的检索增强系统通常包含五个环节。索引环节把语料切分为片段，用编码器映射为向量，并建立近似最近邻索引以支持毫秒级检索，层次化可导航小世界图是常用的索引结构\upcite{malkov2020efficient}。检索环节从查询出发召回候选片段，稀疏检索以词项匹配为基础，稠密检索用双编码器把查询与片段映射到同一向量空间\upcite{karpukhin2020dense}，其编码方式与训练策略的系统梳理见文献\upcite{zhao2022dense}，而后期交互方法在保留词级表示的同时以最大相似度聚合，在效果与效率之间取得折中\upcite{khattab2020colbert}。重排环节用交叉编码器对候选做精细打分，代价是额外的推理开销\upcite{nogueira2019passage}。生成环节决定证据如何进入模型：拼接式直接把片段放入提示，融合式则对每个片段单独编码后在解码器中聚合，从而支持更多片段\upcite{izacard2020leveraging}。评估环节度量检索质量、答案正确性与证据忠实性三类目标，其口径差异是结论不一致的重要来源\upcite{gao2023enabling,es2023ragas}。图~1 给出各环节的衔接关系、由证据充分性判断触发的迭代检索回路，以及第 3 节四个阶段分别关注的部分。

\begin{figure}[H]
\centerline{\includegraphics[width=\linewidth]{rag-framework.pdf}}
{\zihao{5}图1\quad 检索增强生成系统的参考架构。该架构左边为查询与带证据引用的答案，右边为索引构建、检索、证据充分性判断与生成的流水线（据文献\upcite{lewis2020retrieval,gao2023retrieval,karpukhin2020dense,jeong2024adaptive,yu2023generate}整理）}
\label{fig1}
\end{figure}


\subsection{\heiti 与早期工作的关系}

检索增强并非始于大语言模型。开放域问答长期采用"检索器—阅读器"两段式结构，双稠密检索把这一结构推向端到端可训练\upcite{karpukhin2020dense}。在语言建模一侧，最近邻语言模型在token级别用外部数据存储插值输出分布\upcite{khandelwal2020generalization}，REALM\upcite{guu2020realm} 把检索器引入预训练阶段并异步刷新索引，RETRO\upcite{borgeaud2021improving} 则在数万亿token规模的语料上按块检索并通过交叉注意力注入证据。这些工作确立了"非参数化记忆"的基本思想，也暴露出两个后续被反复讨论的问题：检索器与生成器的训练目标不一致，以及索引与模型规模同步增长带来的工程复杂度。

\section{\heiti 技术演进}

已有综述多把检索增强生成划分为朴素、进阶与模块化三个范式\upcite{gao2023retrieval,fan2024survey}；面向知识视角的近期综述虽补充了知识源类型这一维度，划分依据仍是关于“优化模块插在何处”上\upcite{cheng2025survey}。这一划分虽然抓住了"在哪里插入优化模块"这一维度，但是却难以安放近两年出现的两类工作：一类把检索决策交给模型自身，由模型在推理过程中决定何时检索、检索几次\upcite{yao2023react,jin2025search}；另一类改变了知识的组织形态，把平坦的片段集合升级为实体关系图或层次树\upcite{edge2024from,guo2024lightrag}。本文因此把范式扩展为四个阶段，并以"检索与生成的耦合方式"作为划分判据，详见表~1；四个阶段的时间分布与代表性工作如图~2 所示。

\begin{figure}[H]
\centerline{\includegraphics[width=\linewidth]{rag-evolution.pdf}}
{\zihao{5}图2\quad 2020—2025 年代表性工作的演进分布（据本文纳入文献\upcite{lewis2020retrieval,guu2020realm,jin2025search,edge2024from}整理）}
\label{fig2}
\end{figure}


\subsection{\heiti 朴素检索增强}

朴素阶段的共同特征是检索与生成以固定流程耦合：给定查询，先检索固定数量的片段，再一次性生成答案，训练目标为端到端的条件似然。Lewis 等人\upcite{lewis2020retrieval}把检索器视为隐变量，用式~\eqref{eq:rag} 的边缘化目标联合训练检索器与生成器，并给出两种解码方式，其中序列级方式对每个候选文档生成完整答案后加权，token级方式在每一步重新加权，前者适合需要整体一致性的长答案，后者适合短语级答案。REALM\upcite{guu2020realm} 把这一思想前移到预训练阶段，用掩码语言建模信号监督检索器，并通过异步刷新避免每次梯度更新都重建索引。FiD等人\upcite{izacard2020leveraging}改变了证据进入模型的方式：每个片段单独编码，解码器在交叉注意力中聚合多个片段的表示，使片段数量不再受输入长度限制。RETRO等人\upcite{borgeaud2021improving}则在数万亿token的语料上按块检索，用交叉注意力把检索到的邻域信息注入每一层，从而在不显著增加参数量的前提下提升语言建模能力。token级的最近邻语言模型可以看作这一思想的最细粒度版本，它以外部数据存储上的近邻分布对输出做插值\upcite{khandelwal2020generalization}。

朴素阶段确立了检索增强的基本有效性原因，也暴露出三个结构性局限。第一，检索器只在"片段能否提升生成概率"这一间接信号下学习，训练信号稀疏，容易退化为与查询词面相似的检索；第二，索引被视为静态资源，语料更新需要重建索引并重新训练，与"知识可动态更新"的初衷相悖；第三，固定数量的片段与固定的检索时机对问题难度不敏感，面对单个事实型问题足够解决，对需要多步组合或多文档聚合的问题则明显不足\upcite{chen2024benchmarking, trivedi2023interleaving}。这些局限直接催生了进阶阶段。

\subsection{\heiti 进阶检索增强}

进阶阶段的划分依据是：在"检索—生成"的基本骨架两侧插入优化模块，但不改变骨架本身。按插入位置可分为查询侧、召回侧、证据侧与调度侧四类改进。

查询侧改进试图弥合查询与文档之间的语义鸿沟。HyDE\upcite{gao2022precise} 先让语言模型生成一段假设性答案，再以该答案的向量去检索，用生成文本替代短查询作为检索入口；Query2Doc\upcite{wang2023query} 用伪文档扩展查询，在数据稀缺的场景下提升召回；另一类工作把查询改写建模为可训练的策略，用强化学习优化"改写—检索—生成"的整体效果\upcite{ma2023query}。这类方法的共同代价是引入了额外的生成调用，其效果高度依赖生成器对领域知识的掌握程度，在生成器本身缺乏相关知识的场景中，假设性文本会把检索引向错误方向。与查询改写方向相反，另一类工作直接用生成替代检索\upcite{yu2023generate}，先让语言模型生成上下文再据此作答，以生成能力换取证据覆盖。

召回方面改进关注候选集合的质量与多样性。后期交互方法保留词级表示，用最大相似度聚合查询与文档的匹配，在效果与效率之间取得折中\upcite{khattab2020colbert}；稀疏与稠密检索的混合则利用二者互补性，词项匹配擅长精确命名实体，稠密表示擅长语义改写。值得注意的是，零样本检索基准上的系统比较显示，经过调优的词项匹配方法仍是难以超越的强基线，稠密检索并非在所有领域都占优\upcite{thakur2021beir}。这说明"更强的编码器"并不等于"更适合当前语料的检索器"，二者不能相互替代。

依据方面改进关注片段进入模型之前的处理。交叉编码器重排\upcite{nogueira2019passage}用查询与片段的联合编码做精细打分，能够显著提升前若干位的精度，代价是每个候选一次前向计算；上下文压缩\upcite{jiang2023llmlingua}则从另一个方向降低成本，粗到细的提示压缩用小型语言模型逐段删除不重要token，抽取式与生成式压缩\upcite{xu2024recomp}则把多个片段合并为更短的摘要。压缩的风险在于，被删去的往往是连接性表述或限定条件，而这类信息恰恰决定答案是否可信。

调度方面改进决定"何时检索、检索几次"。FLARE \upcite{jiang2023active}在生成过程中监测token概率，当模型置信度下降时才触发检索，把检索时机从"生成之前"改为"生成之中"；IRCoT\upcite{trivedi2023interleaving}把检索与思维链交错，用已生成的推理步骤作为下一轮查询，使多跳问题的每一跳都能获得针对性证据；Self-RAG \upcite{asai2023self}通过反思标记让模型自行判断是否需要检索、证据是否相关、输出是否被证据支持；Adaptive-RAG\upcite{jeong2024adaptive} 则先用一个分类器判断问题复杂度，再在"不检索""单步检索""多步检索"之间路由；Corrective RAG\upcite{yan2024corrective} 引入轻量检索评估器，在证据质量不足时回退到网络检索并对证据做分解再组合。

进阶阶段显著拓宽了 RAG 的适用面，但也带来新的问题：模块越多，延迟与工程复杂度越高；各模块独立优化，其误差沿流程累积；多数模块的超参数（片段数、压缩率、触发阈值）缺乏可迁移的设定依据。这些问题在模块化阶段被显式地当作设计对象处理。

\subsection{\heiti 模块化检索增强}

模块化阶段的判据是：把检索增强拆分为若干可替换、可编排的模块，并研究模块之间的组合规律与训练信号的对齐方式。代表性工作可归为两条线索。

第一条线索关注检索器与生成器的联合优化。Atlas\upcite{izacard2022atlas} 以 FiD 为骨架，在少量样本上联合训练检索器与生成器，表明检索增强模型可以在少样本设定下取得与更大规模参数模型相当的效果；REPLUG\upcite{shi2023replug} 把语言模型视为黑盒，用语言模型对检索文档的打分构造监督信号来训练检索器，使检索器与下游生成目标对齐；RA-DIT\upcite{lin2024dit} 则从两个方向同时做指令微调，既让生成器学会使用检索结果，也让检索器学会面向生成任务排序。三者共同说明，检索器与生成器的目标不一致是朴素阶段遗留的核心问题，而解决它不一定需要端到端联合训练，用生成器的反馈构造检索监督同样可行。

第二条线索关注知识组织形态与编排抽象。RAPTOR\upcite{sarthi2024raptor} 对片段做递归聚类与摘要，构建层次树，使检索可以在不同抽象粒度上进行；MemoRAG\upcite{qian2025memorag} 引入全局记忆模块，先把长文档压缩为记忆表示，再据此生成线索以引导检索；DSP\upcite{khattab2022demonstrate} 把检索与生成抽象为声明式流水线，允许用少量示例编排多阶段查询；面向效率的工作\upcite{jin2024ragcache}则从系统侧切入，通过缓存检索增强场景中的中间状态降低首字延迟。

模块化阶段的价值在于把"怎么组合"变成可研究的问题，其代价是组合空间迅速膨胀：模块顺序、片段数量、是否迭代、是否重排的不同选择构成庞大的配置空间，而现有基准很少报告完整配置，导致复现与比较困难。这也是第 5 节所述评测问题的直接来源。

\subsection{\heiti 智能体化与图结构增强}

最近的工作在两个方向上突破了前三阶段的隐含假设。

第一个方向是检索决策的自主化。ReAct\upcite{yao2023react} 让模型在推理与行动之间交替，行动可以是检索，模型据此在多跳问题上自行组织证据链；Toolformer\upcite{schick2023toolformer} 以自监督方式让模型学会何时调用外部接口；WebGPT\upcite{nakano2021webgpt} 以浏览器操作为行动空间，用人类反馈训练模型检索并引用网页；Search-R1\upcite{jin2025search} 用强化学习直接优化"推理—检索—再推理"的轨迹，以最终答案的正确性作为奖励，使检索时机与查询构造由策略而非启发式规则决定。综述性工作\upcite{singh2025agentic}把这类方法统称为智能体化检索增强，并强调其与工具调用、记忆机制的耦合。

第二个方向是知识组织的结构化。GraphRAG\upcite{edge2024from} 先用语言模型从语料中抽取实体与关系，构建图并对图做社区划分，为每个社区生成摘要；回答全局性问题时以社区摘要为证据，回答局部性问题时沿实体邻域扩展。LightRAG \upcite{guo2024lightrag}采用双层检索，分别面向具体实体与抽象主题，并支持增量更新以降低重构成本。HippoRAG\upcite{gutirrez2024hipporag} 借用海马体索引的类比，用知识图与个性化 PageRank 完成多跳联想，把"一次检索"扩展为"沿图扩散"。近期综述\upcite{zhang2025graphrag}把这一路线归纳为图结构化知识表示、保上下文的图上检索与结构感知的知识融合三个环节，并指出其瓶颈集中在专业领域的复杂查询理解与跨源知识整合。与平坦片段相比，图结构在需要跨文档聚合与多跳推理的问题上更有优势，代价是构图阶段的抽取误差会沿图传播，且图更新与检索的复杂度都显著高于向量索引。

这一阶段还需要回答一个竞争性问题：当上下文窗口扩展到足以容纳整篇文档时，检索是否仍然必要。控制实验显示，模型对长上下文中间位置的信息利用明显弱于首尾位置\upcite{liu2024lost}；而在资源充足的设定下，长上下文模型在部分任务上可以超过检索增强模型，二者结合的混合路由策略则能同时兼顾成本与效果\upcite{li2024retrieval}。近期工作进一步表明，把长文档整体切块后按长上下文方式检索，可以在不增加检索次数的情况下提升证据的完整性\upcite{jiang2024longrag}。这说明检索的价值不仅在于"取回证据"，还在于"筛选证据"，后者在上下文变长后反而更加重要。

\subsection{\heiti 横向比较与适用条件}

表~1 从核心机制、主要优点与适用条件三个方面比较了代表性方法。四个阶段并非彼此替代：朴素阶段针对证据来源，进阶阶段针对证据与问题的匹配，模块化阶段针对模块之间的组合与对齐，智能体与图结构阶段针对检索时机与检索所依据的结构。同一系统常同时包含多个阶段的组件，因此阶段划分用于定位瓶颈，而不是为系统归类。

\end{multicols}

\renewcommand{\arraystretch}{1.25}
\begin{table}[H]
\zihao{5-}
\centering {\heiti 表1\quad 代表性检索增强方法的机制与适用条件}
\vspace {-2.5mm}
\begin{center}
\begin{tabular}{p{22mm}p{27mm}p{50mm}p{55mm}}
\toprule
  方法 & 阶段 & 核心机制 & 适用条件与缺点\\
\hline
RAG & 朴素 & 检索器为隐变量，按式~\eqref{eq:rag} 边缘化 top-$k$ 文档 & 单跳事实型问答；需重训练，索引更新代价高 \\
REALM & 朴素 & 预训练期引入检索器，异步刷新索引 & 知识需持续更新；训练不稳定 \\
FiD & 朴素 & 逐片段编码，解码器交叉注意力聚合 & 证据分散于多段；编码开销随片段数增长 \\
RETRO & 朴素 & 块级检索，交叉注意力注入每一层 & 万亿级语料；需从零训练 \\
kNN-LM & 朴素 & token级最近邻分布插值 & 无需改结构；存储与延迟开销大 \\
HyDE & 进阶 & 生成假设性答案后再检索 & 查询短、语义鸿沟大；依赖生成质量 \\
Query2Doc & 进阶 & 伪文档扩展查询 & 标注稀缺领域；可能引入生成幻觉 \\
LLMLingua & 进阶 & 粗到细的提示压缩 & 上下文超长；压缩可能丢失限定条件 \\
RECOMP & 进阶 & 抽取式或生成式证据压缩 & 多片段冗余；摘要错误影响忠实性 \\
FLARE & 进阶 & 置信度下降时触发检索 & 长答案生成；触发阈值敏感 \\
IRCoT & 进阶 & 检索与思维链交错 & 多跳推理；调用次数增加 \\
Self-RAG & 进阶 & 反思标记控制检索与自我批判 & 需专门训练；标注成本高 \\
Adaptive-RAG & 进阶 & 按问题复杂度路由检索策略 & 查询难度分布不均；分类器误差传导 \\
Atlas & 模块化 & FiD 骨架上的少样本联合训练 & 少样本设定；训练资源需求高 \\
REPLUG & 模块化 & 用语言模型反馈训练检索器 & 黑盒生成器；需额外打分开销 \\
RA-DIT & 模块化 & 检索器与生成器双向指令微调 & 需要监督微调数据 \\
RAPTOR & 模块化 & 递归聚类摘要构成层次树 & 需多粒度证据；摘要质量决定上限 \\
GraphRAG & 智能体/图 & 实体关系图与社区摘要 & 全局性问题；构图成本高、更新困难 \\
LightRAG & 智能体/图 & 双层检索与增量图更新 & 动态语料；图质量依赖抽取 \\
HippoRAG & 智能体/图 & 知识图与个性化 PageRank 扩散 & 多跳联想；图规模受限 \\
Search-R1 & 智能体/图 & 强化学习优化检索轨迹 & 需可交互检索环境与奖励设计 \\
\bottomrule
\end{tabular}
\label{tab1}
\end{center}
\vspace{-1mm}
{\zihao{6}\centering 注：表中"阶段"按本文第 3 节的四阶段划分；机制与条件据文献\upcite{lewis2020retrieval,guu2020realm,izacard2020leveraging,borgeaud2021improving,khandelwal2020generalization,gao2022precise,wang2023query,jiang2023llmlingua,xu2024recomp,jiang2023active,trivedi2023interleaving,asai2023self,jeong2024adaptive,izacard2022atlas,shi2023replug,lin2024dit,sarthi2024raptor,edge2024from,guo2024lightrag,gutirrez2024hipporag,jin2025search}整理。}
\end{table}

\begin{multicols}{2}

\section{\heiti 评测体系与实证结论}

\subsection{\heiti 基准数据集与任务形态}

评测的演进与方法的演进同步。早期工作沿用开放域问答的基准，以精确匹配与token级 F1 为主指标；KILT 把五个知识密集型任务统一到同一接口下，并要求系统给出证据出处，使"答案正确"与"证据正确"第一次被同时度量\upcite{petroni2021kilt}。BEIR 则从检索一侧出发，用十八个跨领域数据集检验零样本迁移能力，其结论对大语言模型时代的检索器设计仍有约束力\upcite{thakur2021beir}。

针对检索增强特有的失效模式，后续基准在设计上更有针对性。HotpotQA、2WikiMultiHopQA 与 MuSiQue 通过构造多跳问题，要求系统跨文档组合证据\upcite{yang2018hotpotqa,ho2020constructing,trivedi2022musique}；RGB 显式区分噪声鲁棒性、负样本拒答、信息整合与反事实鲁棒性四类能力，把"检索到无关或矛盾证据时会发生什么"变成可测项\upcite{chen2024benchmarking}；CRAG 覆盖多个领域与多种问题类型，并同时提供检索与生成的评测口径\upcite{yang2024crag}；MultiHop-RAG 面向多跳查询构造语料与标注\upcite{tang2024multihop}。在可信性一侧，ALCE 要求模型为每一句结论给出引用并度量引用的正确性与完整性\upcite{gao2023enabling}，RAGTruth 则以人工标注的幻觉语料刻画检索增强场景下的三类幻觉\upcite{niu2023ragtruth}。表~2 汇总了常用基准的考察重点。

\end{multicols}

\renewcommand{\arraystretch}{1.2}
\begin{table}[H]
\zihao{5-}
\centering {\heiti 表2\quad 常用基准与主要考察点}
\vspace {-2.5mm}
\begin{center}
\begin{tabular}{p{26mm}p{34mm}p{88mm}}
\toprule
基准 & 任务形态 & 主要考察点 \\
\hline
KILT & 五类知识密集型任务 & 答案正确性与证据出处的一致性 \\
BEIR & 跨领域零样本检索 & 检索器的领域迁移能力 \\
HotpotQA & 多跳问答 & 跨文档证据组合 \\
2WikiMultiHopQA & 多跳问答 & 推理步序与证据可解释性 \\
MuSiQue & 组合式多跳问答 & 干扰项下的组合推理 \\
RGB & 检索增强问答 & 噪声鲁棒性、拒答、信息整合、反事实鲁棒性 \\
CRAG & 多领域问答 & 检索与生成的分层评测 \\
MultiHop-RAG & 多跳查询 & 多文档多跳检索 \\
ALCE & 带引用的生成 & 引用的正确性与完整性 \\
RAGTruth & 幻觉标注语料 & 检索增强场景的幻觉类型与强度 \\
\bottomrule
\end{tabular}
\label{tab2}
\end{center}
\vspace{-1mm}
{\zihao{6}\centering 注：考察点据文献\upcite{petroni2021kilt,thakur2021beir,yang2018hotpotqa,ho2020constructing,trivedi2022musique,chen2024benchmarking,yang2024crag,tang2024multihop,gao2023enabling,niu2023ragtruth}整理。}
\end{table}

\begin{multicols}{2}

\subsection{\heiti 评价指标}

已有综述把评测要素拆解为评测对象、数据集与指标三个维度，并把指标进一步区分为检索质量、答案忠实性与答案相关性\upcite{yu2024evaluation}。据此，本文把指标分为三层。检索层沿用召回率与归一化折损累计增益，衡量正确证据是否进入候选集以及排序是否合理；生成层使用精确匹配与 F1 衡量答案字面一致性，但在生成式答案上这两项指标与人工判断的相关性有限，因此越来越多工作引入基于自然语言推理的忠实性判定，检验答案的每个论断是否能被给定证据支持。系统层则关注成本，包括首字延迟、总时延、检索调用次数与token消耗，这一层在过去常被忽略，却是判断方法能否落地的前提。

为降低人工标注成本，RAGAS 提出一组无需参考答案的指标，用大语言模型分别估计忠实性、答案相关性与上下文相关性\upcite{es2023ragas}；ARES 用少量标注数据微调判别器，并结合预测驱动推断给出带置信区间的评测结果\upcite{saadfalcon2024ares}；RAGChecker 进一步把评测细化到论断级别，分别诊断检索器与生成器的贡献\upcite{ru2024ragchecker}。这类"以模型评模型"的做法提高了效率，但也引入了评判者偏差：当评判模型与被评模型同源时，评测结果可能系统性偏高，目前尚缺少对这一偏差的定量刻画。

\subsection{\heiti 共识与分歧}

综合现有证据，可以得到三点较为稳固的共识。第一，证据质量决定系统上限。多篇工作在受控条件下替换检索器而不改动生成器，性能随之显著变化，说明检索环节是当前的主要瓶颈\upcite{chen2024benchmarking,yang2024crag}。第二，证据的位置与噪声不可忽略。长上下文实验显示模型对中间位置信息的利用率明显低于首尾\upcite{liu2024lost}，而噪声鲁棒性实验显示无关片段的比例上升会持续削弱答案质量\upcite{chen2024benchmarking}。第三，简单而强的基线不容忽视。跨领域检索评测表明调优后的词项匹配仍具竞争力\upcite{thakur2021beir}，黑盒拼接式增强在多个任务上也能取得稳定效果\upcite{ram2023context}，这提示复杂流水线的增益需要与同等成本的基线对比才有意义。

分歧主要集中在三处。其一，检索增强与长上下文的关系。有研究在资源充足时发现长上下文模型可以超过检索增强模型，并提出按需路由的混合方案\upcite{li2024retrieval}；但也有工作指出长上下文模型在证据分散时仍会遗漏关键信息，且推理成本随长度快速上升\upcite{liu2024lost}。两者的结论并不真正矛盾，差异来自成本约束是否被纳入比较。其二，联合训练与解耦训练的取舍。Atlas 主张端到端联合训练\upcite{izacard2022atlas}，REPLUG 与 RA-DIT 则证明以生成器反馈监督检索器同样有效\upcite{shi2023replug,lin2024dit}，但三者在训练数据、模型规模与评测设定上均不相同，缺少同设定下的对照实验。其三，迭代检索的必要性。IRCoT 与 FLARE 在多跳与长答案任务上效果明显\upcite{trivedi2023interleaving,jiang2023active}，而在单跳任务上多次检索带来的增益常被额外延迟抵消，目前尚无统一结论说明何时值得迭代。

此外还存在明显的证据缺口：图结构方法的效果究竟来自结构本身，还是来自更多的上下文与更多的模型调用，现有工作报告的多为端到端指标而缺少针对性消融\upcite{edge2024from,guo2024lightrag}；中文与专业领域的基准相对缺乏，跨语言迁移的系统性证据不足；不少方法未公开完整配置，复现成本高，削弱了结论的可比性。

\section{\heiti 问题与展望}

\subsection{\heiti 现存问题}

第一，检索器的训练信号仍然薄弱。无论是用生成概率还是用最终答案正确性作为监督，信号都要经过生成器才能回传，信噪比低；用大语言模型打分替代虽缓解了这一问题\upcite{shi2023replug}，却把生成器的偏差引入了检索器。

第二，噪声与矛盾的鲁棒性不足。现有系统在检索到无关片段时倾向于"强行使用"，在检索到与参数知识冲突的片段时则表现出两种极端：要么盲从片段，要么固守参数知识。RGB 与 RAGTruth 的结果表明，这两种失效都与证据的组织方式有关，而非单纯的模型能力问题\upcite{chen2024benchmarking,niu2023ragtruth}。

第三，多跳与全局性问题仍未解决。单次检索加单次生成的骨架无法完成需要跨文档聚合的问题，迭代检索与图结构是两条可行路径，但前者调用次数不可控，后者构图成本高且抽取误差会沿图传播\upcite{trivedi2023interleaving,edge2024from}。

第四，效率与成本缺少统一口径。缓存、压缩与稀疏化各自报告了加速比，但测试硬件、批量大小与基线实现差异很大，跨工作的比较意义有限\upcite{jin2024ragcache,jiang2023llmlingua,xu2024recomp}。

第五，安全与可信问题开始显现。知识库是新的攻击面，少量投毒文档即可显著改变模型输出\upcite{zou2025poisonedrag}；而可溯源性方面，即使系统给出了引用，引用的正确性与完整性也远未达到可用水平\upcite{gao2023enabling,liu2023evaluating}。

\subsection{\heiti 未来研究方向}

其一，代价感知的自适应检索。现有自适应方法多以问题复杂度为路由依据\upcite{jeong2024adaptive}，未来需要把延迟预算、检索费用与置信度纳入同一决策框架，让系统在给定成本下最大化证据效果，而不是在固定预算下追求单点指标。

其二，端到端联合优化与可微分检索。检索的离散性使梯度难以穿透，现有替代方案包括强化学习\upcite{jin2025search}与基于生成器反馈的监督\upcite{shi2023replug,lin2024dit}。如何在保持索引可更新的前提下实现稳定的联合训练，仍是有价值的问题。

其三，结构化与多模态知识的接入。图结构在聚合性问题上展现了潜力\upcite{edge2024from,guo2024lightrag,gutirrez2024hipporag}，但构图、更新与检索三者尚未形成统一框架；多模态证据的检索与融合也缺少与文本证据可比的组织方式。

其四，可溯源生成与引用质量。ALCE 把引用质量拆分为正确性与完整性两个维度\upcite{gao2023enabling}，后续工作需要在训练阶段就引入引用约束，而非在生成后附加引用。对生成式搜索引擎的评估显示，现有系统的引用支持率仍有较大提升空间\upcite{liu2023evaluating}。

其五，统一评测协议。建议后续工作在报告效果的同时报告检索调用次数、检索延迟与token消耗，并在同一实现下给出简单基线与复杂方法的对照，这是缓解当前结论冲突最直接的办法\upcite{ru2024ragchecker}。

\section{\heiti 结论}

本文以技术演进为主线梳理了大语言模型检索增强生成的研究进展，将其划分为朴素检索增强、进阶检索增强、模块化检索增强以及智能体化与图结构增强四个阶段，并给出各阶段的核心机制、代表性方法与适用条件。分析表明，四个阶段解决的是不同层面的问题：证据从何而来、证据是否对题、模块如何组合、何时并沿何种结构检索；它们之间是叠加而非替代关系。

对评测证据的梳理显示，领域内已就"证据质量决定上限""位置与噪声不可忽略""简单强基线不容忽视"三点形成共识，而在检索增强与长上下文的关系、联合训练与解耦训练的取舍、迭代检索的必要性三处仍存在分歧，分歧多源于成本口径与实验设定的不一致。据此，本文提出代价感知的自适应检索、端到端联合优化、结构化与多模态知识接入、可溯源生成以及统一评测协议五个方向。

归根结底，检索增强生成的效果不取决于模块数量，而取决于检索时机、证据粒度与生成器利用证据的方式三者是否匹配。未来的系统设计宜以这一匹配关系为中心，而非继续堆叠独立的优化模块。

\vspace {3mm}
\centerline{\zihao{5}\textsf{参~考~文~献}}
\zihao{5-}
\bibliographystyle{gbt7714-numerical}
\bibliography{rag.bib}

\end{multicols}
\end{document}
